\documentclass[11pt,a4paper]{article}

\usepackage[a4paper,margin=1in]{geometry}
\usepackage{fontspec}
\IfFontExistsTF{Times New Roman}
  {\setmainfont{Times New Roman}}
  {\IfFontExistsTF{TeX Gyre Termes}{\setmainfont{TeX Gyre Termes}}{\setmainfont{Times}}}
\usepackage{newtxmath}
\usepackage{amsmath}
\usepackage{booktabs}
\usepackage{longtable}
\usepackage{array}
\usepackage{enumitem}
\usepackage{xcolor}
\usepackage{listings}
\usepackage{graphicx}
\usepackage{float}
\usepackage[hidelinks]{hyperref}
\usepackage{xurl}
\usepackage{fancyhdr}
\usepackage{microtype}

\setlength{\parindent}{0pt}
\setlength{\parskip}{0.65em}
\setlength{\emergencystretch}{3em}
\setlist[itemize]{leftmargin=1.4em,itemsep=0.2em,topsep=0.2em}
\setlist[enumerate]{leftmargin=1.6em,itemsep=0.2em,topsep=0.2em}

\pagestyle{fancy}
\fancyhf{}
\lhead{\small Ticket 08 Study Report}
\rhead{\small Obstacle and Complementarity Diagnostics}
\cfoot{\thepage}
\renewcommand{\headrulewidth}{0.3pt}

\lstset{
  basicstyle=\ttfamily\small,
  breaklines=true,
  frame=single,
  columns=fullflexible,
  keepspaces=true,
  rulecolor=\color{black!30},
  xleftmargin=0.5em,
  xrightmargin=0.5em
}

\newcommand{\code}[1]{\texttt{#1}}
\newcommand{\file}[1]{\path{#1}}
\newcommand{\func}[1]{\path{#1}}
\newcolumntype{L}[1]{>{\raggedright\arraybackslash}p{#1}}
\newcolumntype{P}[1]{>{\raggedright\arraybackslash\ttfamily\footnotesize}p{#1}}

\title{\textbf{Ticket 08 Study Report: Obstacle and Complementarity Diagnostics}\\
\large SURF2026 American Option Risk Surfaces}
\author{Codex-assisted implementation report}
\date{June 16, 2026}

\begin{document}
\maketitle

\begin{abstract}
Ticket 08 adds a diagnostic layer for American CN/PSOR results. It does not introduce a new
pricing algorithm. Instead, it reconstructs the discrete linear complementarity problem solved at
each American time step and checks the obstacle gap, equation gap, and complementarity product on
interior spot nodes. The diagnostics are applied to the representative American put case from the
validation ladder and the positive-dividend American call case from Ticket 07. This report is
written for beginner researchers with accounting, data, machine-learning, and Python experience
who are still learning obstacle problems, LCP diagnostics, and CN/PSOR validation.
\end{abstract}

\tableofcontents
\newpage

\section{Role of Ticket 08 in the Whole Solver Project}

Tickets 01 through 07 built the pricing and validation ladder: analytic Black-Scholes utilities,
finite-difference grids, European Crank--Nicolson validation, American CN/PSOR solving,
American put validation, no-dividend American call validation, and dividend-paying American call
validation. Ticket 08 adds a diagnostic layer on top of those American solver outputs.

The key purpose is to check more than price levels. A numerical American option result can look
reasonable and still fail the LCP logic. Ticket 08 therefore asks:
\begin{itemize}
  \item Does the value stay above payoff?
  \item Does the reconstructed equation inequality hold?
  \item Is the product of the two LCP gaps close to zero?
  \item Did the PSOR iterations converge while these diagnostics were computed?
\end{itemize}

This ticket is not boundary extraction. It does not locate a free boundary, does not create
Greeks, does not create stress maps, does not create datasets, and does not modify the solver.

\section{Theory and Methodology/Literature Connection}

The documented source basis for Ticket 08 is internal:
\begin{itemize}
  \item \file{reports/01_literature/literature_map.md};
  \item \file{reports/02_math/formulation_note.md};
  \item \file{reports/03_solver/solver_validation_plan.md};
  \item \file{reports/03_solver/solver_implementation_tickets.md};
  \item the Ticket 01 through Ticket 07 implementation reports.
\end{itemize}

The formulation note explains that American exercise turns the pricing problem into a
variational inequality and, after discretization, a linear complementarity problem. The solver
validation plan explicitly calls for obstacle violation and complementarity residual reporting.
The implementation ticket file says the diagnostics must keep boundary nodes separate from
interior LCP residuals and must not hide maximum residuals behind averages.

No BibTeX file is created for Ticket 08. The report uses the internal project documents as the
verified source basis and avoids inventing external bibliographic details.

\section{Obstacle, Equation Gap, and Complementarity in Plain Language}

For an American option, the value \(U\) must be at least the payoff \(\Phi\):
\begin{equation}
  U \geq \Phi.
\end{equation}
This is the payoff obstacle. If a computed American value falls below payoff, the holder would be
better off exercising immediately, so the numerical value is financially invalid.

The CN/PSOR solver also solves a continuation equation at every time step. In matrix form the
interior-node LCP is:
\begin{equation}
  A u \geq b,\qquad u \geq \Phi,\qquad
  (u-\Phi)^\top(Au-b)=0.
\end{equation}

The two gaps have plain meanings:
\begin{itemize}
  \item \(u-\Phi\): how far the value is above immediate exercise;
  \item \(Au-b\): how far the discrete equation inequality is from binding.
\end{itemize}

Complementarity says that each interior node should be in one of two modes. In a continuation-like
node, \(u>\Phi\) and the equation should be close to equality. In an exercise-like node,
\(u=\Phi\) and the equation inequality can be positive. The product of the two gaps should
therefore be close to zero.

\section{Why Diagnostics Come After Put and Call Validation}

Ticket 08 comes after Tickets 05, 06, and 07 because diagnostics need solver outputs to inspect.
Those earlier tickets established representative American put and call cases:
\begin{itemize}
  \item American put: early exercise can matter and the value should dominate payoff and European
        put values.
  \item No-dividend American call: the American value should approximately match the European
        call when \(q=0\).
  \item Dividend-paying American call: positive dividend yield can create meaningful
        American-minus-European value.
\end{itemize}

Ticket 08 does not repeat those price-validation arguments. It checks whether the American
solver outputs also satisfy the internal LCP conditions expected from the CN/PSOR formulation.

\section{LCP Diagnostic Formulas}

For each time step and each interior spot node, Ticket 08 computes:
\begin{align}
  \text{value gap} &= u - \Phi,\\
  \text{equation gap} &= Au - b,\\
  \text{obstacle violation} &= \max(\Phi-u,0),\\
  \text{equation violation} &= \max(b-Au,0),\\
  \text{complementarity product} &= (u-\Phi)(Au-b).
\end{align}

\begin{longtable}{L{0.28\textwidth}L{0.60\textwidth}}
\toprule
Quantity & Beginner meaning \\
\midrule
\endhead
Value gap & Positive means the option value is above immediate exercise. Zero means the value is
at the payoff obstacle. \\
Equation gap & Positive means the equation inequality has slack. Near zero means the continuation
equation is binding. \\
Obstacle violation & Any positive value means the computed option value fell below payoff. \\
Equation violation & Any positive value means the reconstructed matrix inequality \(Au\geq b\)
was violated. \\
Complementarity product & Should be close to zero. A large value means the node is neither cleanly
continuation-like nor exercise-like under the LCP conditions. \\
\bottomrule
\end{longtable}

\section{Diagnostics Versus Boundary Extraction}

Ticket 08 deliberately avoids boundary extraction. It computes value gaps and diagnostic node
counts, but it does not search for a transition location, does not interpolate an exercise
boundary, and does not label any figure with a boundary. Terms such as exercise-like and
continuation-like are local diagnostic classifications only.

This distinction matters. Diagnostics answer: ``does the computed value satisfy the LCP
conditions?'' Boundary extraction asks: ``where is the transition curve?'' The second question
requires a separate algorithm and belongs to a later ticket.

\section{Implementation Design Decisions}

\begin{longtable}{L{0.30\textwidth}L{0.58\textwidth}}
\toprule
Decision & Reason \\
\midrule
\endhead
Create \file{src/american_risk_surfaces/diagnostics/lcp.py} &
Diagnostics are separate from the solver core, keeping pricing logic and validation logic
independent. \\
Use interior nodes only &
Boundary nodes are imposed values, not unknowns solved by the LCP. Including them would pollute
the residual. \\
Reconstruct the Ticket 04 matrix form &
The diagnostic must check the same \(A\) and \(b\) that the solver used, including new-time
boundary adjustments. \\
Report maxima and minima &
Maximum residuals reveal localized failures that averages can hide. \\
Use representative put and dividend-call cases &
The ticket checks diagnostics on two important American regimes without becoming a full grid
sensitivity study. \\
Use three simple figures &
Obstacle violation, equation violation, and complementarity product each answer a different
diagnostic question. \\
\bottomrule
\end{longtable}

\section{Code-Level Function Explanations}

\begin{longtable}{L{0.38\textwidth}L{0.50\textwidth}}
\toprule
Function or class & Role \\
\midrule
\endhead
\func{LCPGapArrays} & Stores the five interior-node gap arrays for one time step. \\
\func{ReconstructedLCPStep} & Stores \(u\), \(\Phi\), \(Au\), \(b\), and gaps for one time step. \\
\func{StepDiagnosticRow} & Stores scalar metrics for one time step, including PSOR metadata. \\
\func{LCPDiagnosticSummary} & Stores whole-result maxima, tolerances, convergence status, and
diagnostic status. \\
\func{compute_lcp_gap_arrays} & Computes value gap, equation gap, violations, and complementarity
product from arrays. \\
\func{reconstruct_lcp_step} & Rebuilds the interior-node LCP for one stored American CN/PSOR step. \\
\func{summarize_lcp_step} & Converts reconstructed arrays into scalar diagnostics and node counts. \\
\func{diagnose_lcp_result} & Runs reconstruction and summary over every PSOR time step. \\
\func{run_diagnostic_cases} & Runs the representative experiment cases and prepares CSV/figure
metadata. \\
\bottomrule
\end{longtable}

\section{Diagnostic Cases and Parameters}

\begin{longtable}{@{}L{0.16\textwidth}L{0.10\textwidth}L{0.07\textwidth}L{0.07\textwidth}L{0.07\textwidth}L{0.07\textwidth}L{0.09\textwidth}L{0.07\textwidth}L{0.07\textwidth}@{}}
\toprule
Case & Option & \(K\) & \(T\) & \(r\) & \(q\) & \(\sigma\) & \(M\) & \(N\) \\
\midrule
\endhead
Put medium & put & 1.0 & 1.0 & 0.05 & 0.02 & 0.20 & 120 & 120 \\
Dividend call & call & 1.0 & 1.0 & 0.05 & 0.08 & 0.20 & 120 & 120 \\
\bottomrule
\end{longtable}

Both cases use \(S_{\max}=4.0\). The purpose is diagnostic correctness and interpretation, not
full grid or domain sensitivity.

\section{Diagnostic Result Tables and Interpretation}

\subsection{Summary Diagnostics}

{\footnotesize
\setlength{\tabcolsep}{3pt}
\begin{longtable}{@{}L{0.12\textwidth}L{0.07\textwidth}L{0.10\textwidth}L{0.10\textwidth}L{0.11\textwidth}L{0.11\textwidth}L{0.08\textwidth}@{}}
\toprule
Case & Status & Max obstacle & Max equation & Max comp. product & Max final update & Max iter. \\
\midrule
\endhead
Put & PASS & 0 & \(1.5566\times10^{-9}\) & \(6.6483\times10^{-11}\) & \(9.9971\times10^{-9}\) & 10 \\
Call & PASS & 0 & \(2.1442\times10^{-9}\) & \(2.3796\times10^{-11}\) & \(9.4309\times10^{-9}\) & 10 \\
\bottomrule
\end{longtable}
}

Both cases pass the Ticket 08 diagnostic thresholds. The maximum obstacle violation is zero in
both cases. Equation violations are about \(10^{-9}\), and complementarity products are about
\(10^{-11}\). These numbers are below the diagnostic thresholds of \(10^{-6}\) for equation
violation and complementarity product.

\subsection{Node Classification Summary}

\begin{longtable}{@{}L{0.24\textwidth}L{0.18\textwidth}L{0.20\textwidth}L{0.18\textwidth}@{}}
\toprule
Case & Max exercise-like nodes & Max continuation-like nodes & Max ambiguous nodes \\
\midrule
\endhead
american\_put\_medium & 113 & 51 & 0 \\
dividend\_call\_medium & 114 & 25 & 0 \\
\bottomrule
\end{longtable}

The node classifications are diagnostic labels only. They show that nodes can be separated into
exercise-like and continuation-like categories under the stated tolerances. They do not locate a
free boundary.

\subsection{Selected By-Step Rows}

{\footnotesize
\setlength{\tabcolsep}{3pt}
\begin{longtable}{@{}L{0.18\textwidth}L{0.06\textwidth}L{0.08\textwidth}L{0.08\textwidth}L{0.10\textwidth}L{0.10\textwidth}L{0.10\textwidth}L{0.08\textwidth}L{0.08\textwidth}@{}}
\toprule
Case & Step & \(\tau\) & Iter. & Final update & Max eq. viol. & Max comp. & Ex-like & Cont-like \\
\midrule
\endhead
put & 1 & 0.00833 & 10 & \(3.3486\times10^{-9}\) & \(3.3576\times10^{-11}\) & \(9.5176\times10^{-13}\) & 113 & 6 \\
put & 60 & 0.50000 & 7 & \(1.7795\times10^{-9}\) & \(1.6431\times10^{-10}\) & \(6.3703\times10^{-12}\) & 86 & 33 \\
put & 120 & 1.00000 & 6 & \(8.8380\times10^{-9}\) & \(1.0091\times10^{-9}\) & \(6.4760\times10^{-11}\) & 68 & 51 \\
call & 1 & 0.00833 & 10 & \(5.8695\times10^{-9}\) & \(1.2219\times10^{-9}\) & \(1.2466\times10^{-12}\) & 114 & 5 \\
call & 60 & 0.50000 & 7 & \(4.5496\times10^{-9}\) & \(4.2995\times10^{-10}\) & \(1.2760\times10^{-11}\) & 98 & 21 \\
call & 120 & 1.00000 & 7 & \(3.8139\times10^{-9}\) & \(1.1096\times10^{-9}\) & \(1.3338\times10^{-11}\) & 94 & 25 \\
\bottomrule
\end{longtable}
}

The selected rows show that diagnostics remain small early, halfway through, and at final
maturity horizon. The full by-step CSV contains all 240 rows.

\section{Figure Interpretation}

Figure~\ref{fig:obstacle} plots maximum obstacle violation by time step. Both cases remain at
zero on the generated grid, showing that the stored American values do not fall below payoff.

\begin{figure}[H]
\centering
\includegraphics[width=0.82\textwidth]{results/01_solver_validation/figures/ticket_08_lcp_obstacle_violation_by_step.png}
\caption{Ticket 08 maximum obstacle violation by time step.}
\label{fig:obstacle}
\end{figure}

Figure~\ref{fig:equation} plots maximum equation violation by time step. The values remain near
the PSOR tolerance scale and well below the Ticket 08 review threshold.

\begin{figure}[H]
\centering
\includegraphics[width=0.82\textwidth]{results/01_solver_validation/figures/ticket_08_lcp_equation_violation_by_step.png}
\caption{Ticket 08 maximum equation violation by time step.}
\label{fig:equation}
\end{figure}

Figure~\ref{fig:complementarity} plots maximum absolute complementarity product by time step.
Small values indicate that nodes are close to the LCP alternatives: either continuation-like or
exercise-like.

\begin{figure}[H]
\centering
\includegraphics[width=0.82\textwidth]{results/01_solver_validation/figures/ticket_08_lcp_complementarity_by_step.png}
\caption{Ticket 08 maximum absolute complementarity product by time step.}
\label{fig:complementarity}
\end{figure}

\section{Testing Strategy}

\begin{longtable}{L{0.28\textwidth}L{0.30\textwidth}L{0.30\textwidth}}
\toprule
Test category & What it proves & What it does not prove \\
\midrule
\endhead
Manual gap arrays & Basic formulas and signs are correct on small arrays. & Full solver
diagnostics. \\
Toy LCP classification & Exercise-like and continuation-like counts behave as intended. & Boundary
location. \\
Interior reconstruction & Boundary nodes are excluded from residual arrays. & Boundary values are
economically perfect. \\
American put diagnostics & Representative put residuals are finite and small. & Full put grid
sensitivity. \\
Dividend-call diagnostics & Representative positive-dividend call residuals are finite and small. &
Full dividend sensitivity. \\
CSV artifact tests & Experiment outputs have stable schemas. & That later reports need no extra
columns. \\
Figure tests & Meaningful diagnostic figures are created when matplotlib is available. & That plots
alone prove correctness. \\
Scope tests & Ticket 08 does not introduce forbidden APIs/files. & Future tickets that intentionally
add those features. \\
\bottomrule
\end{longtable}

\section{Test Results and Interpretation}

The final unittest command was:
\begin{lstlisting}[language=bash]
PYTHONPYCACHEPREFIX=/private/tmp/codex_pycache PYTHONPATH=src MPLCONFIGDIR=/private/tmp/matplotlib-cache FC_CACHEDIR=/private/tmp/fontconfig-cache python3 -m unittest discover -s tests
\end{lstlisting}

Observed result:
\begin{lstlisting}
Ran 87 tests

OK
\end{lstlisting}

The Ticket 08 experiment command was:
\begin{lstlisting}[language=bash]
PYTHONPYCACHEPREFIX=/private/tmp/codex_pycache PYTHONPATH=src MPLCONFIGDIR=/private/tmp/matplotlib-cache FC_CACHEDIR=/private/tmp/fontconfig-cache python3 experiments/05_obstacle_complementarity_diagnostics.py
\end{lstlisting}

It wrote two summary rows, 240 by-step rows, and three PNG figures.

\section{Implementation Pitfalls and Debugging Notes}

\begin{longtable}{@{}L{0.23\textwidth}L{0.24\textwidth}L{0.22\textwidth}L{0.19\textwidth}@{}}
\toprule
Recorded item or pitfall & Symptom & Fix or guard & Beginner lesson \\
\midrule
\endhead
Expected TDD red run & The new tests failed because \file{diagnostics/lcp.py} and the Ticket 08
experiment script did not exist yet. & Add the diagnostics package and experiment script. & A red
test confirms the tests target new behavior. \\
Malformed targeted unittest command & \file{tests.test_lcp_diagnostics} could not import because
\file{tests/} is not a package. & Use unittest discovery with \code{-s tests} and \code{-k}. &
Match test commands to the repository's import layout. \\
Wrong equation-gap sign & Reporting \(b-Au\) as the gap would reverse the LCP inequality. & Define
equation gap as \(Au-b\) and equation violation as \(\max(b-Au,0)\). & Sign conventions must be
written down before interpreting residuals. \\
Including boundary nodes & Boundary values are imposed, not solved by the interior LCP. & Slice all
diagnostic arrays with \code{1:-1}. & Residuals should match the unknowns solved by PSOR. \\
Overinterpreting classifications & Exercise-like counts can look like an exercise region. & State
that counts are diagnostic labels only. & Diagnostics are not boundary extraction. \\
Ignoring PSOR convergence & Small gaps are less meaningful if the LCP solve did not converge. &
Report PSOR iterations, final update, and convergence status. & Solver metadata is part of
diagnostic credibility. \\
\bottomrule
\end{longtable}

\section{Limitations}

\begin{longtable}{L{0.34\textwidth}L{0.54\textwidth}}
\toprule
Not included in Ticket 08 & Reason \\
\midrule
\endhead
New pricing algorithm & The ticket diagnoses existing American CN/PSOR outputs only. \\
Free-boundary extraction & Boundary location belongs to a later ticket. \\
Continuation-premium boundary module & Ticket 08 uses local gaps only, not transition extraction. \\
Greeks & Delta and Gamma diagnostics are later work. \\
Stress maps, datasets, neural networks & These require a more complete validated solver pipeline. \\
Full grid/domain sensitivity & The two medium-grid cases are representative diagnostics, not a
convergence study. \\
Formal proof of complementarity convergence & The report gives numerical diagnostic evidence only. \\
\bottomrule
\end{longtable}

\section{Lessons Learned / Study Notes}

The first lesson is that obstacle validity is necessary but not sufficient. A value can stay
above payoff while still failing the equation side of the LCP.

The second lesson is that complementarity is a local either-or condition. At each interior node,
either the option behaves like a continuation PDE point, or it sits on the payoff obstacle.

The third lesson is that boundary nodes should not be included in LCP residuals. They are imposed
by boundary conditions and are not unknowns solved by PSOR.

The fourth lesson is that diagnostic plots should answer diagnostic questions. Ticket 08 uses
figures for obstacle violation, equation violation, and complementarity product; it avoids
figures that could be mistaken for boundary extraction.

\section{Link to the Next Ticket}

Ticket 08 provides the reusable diagnostic layer needed before boundary extraction. The next
ticket can use value and complementarity evidence with more confidence, but it must still define
its own boundary extraction rules, thresholds, interpolation choices, and reporting language.

\section{Reproducibility Record}

\begin{longtable}{L{0.30\textwidth}L{0.58\textwidth}}
\toprule
Item & Record \\
\midrule
\endhead
Diagnostics package & \file{src/american_risk_surfaces/diagnostics/lcp.py}. \\
Experiment & \file{experiments/05_obstacle_complementarity_diagnostics.py}. \\
Tests & \file{tests/test_lcp_diagnostics.py}. \\
Summary CSV & \file{results/01_solver_validation/tables/ticket_08_lcp_diagnostics_summary.csv}. \\
By-step CSV & \file{results/01_solver_validation/tables/ticket_08_lcp_diagnostics_by_step.csv}. \\
Figures & \file{results/01_solver_validation/figures/ticket_08_lcp_obstacle_violation_by_step.png};
\file{results/01_solver_validation/figures/ticket_08_lcp_equation_violation_by_step.png};
\file{results/01_solver_validation/figures/ticket_08_lcp_complementarity_by_step.png}. \\
Report source & \file{reports/03_solver/tickets/ticket_08_obstacle_complementarity_diagnostics.tex}. \\
Report PDF & \file{reports/03_solver/tickets/ticket_08_obstacle_complementarity_diagnostics.pdf}. \\
\bottomrule
\end{longtable}

Commands used:
\begin{lstlisting}[language=bash]
PYTHONPYCACHEPREFIX=/private/tmp/codex_pycache PYTHONPATH=src MPLCONFIGDIR=/private/tmp/matplotlib-cache FC_CACHEDIR=/private/tmp/fontconfig-cache python3 experiments/05_obstacle_complementarity_diagnostics.py
PYTHONPYCACHEPREFIX=/private/tmp/codex_pycache PYTHONPATH=src MPLCONFIGDIR=/private/tmp/matplotlib-cache FC_CACHEDIR=/private/tmp/fontconfig-cache python3 -m unittest discover -s tests
PYTHONPYCACHEPREFIX=/private/tmp/codex_pycache PYTHONPATH=src python3 -m compileall -q src tests experiments
\end{lstlisting}

\end{document}
